%======================================================================
%  style.tex -- page layout and typography
%
%  Same visual language as the "Nonlinear Problems in Mechanics" lecture
%  notes (the RG style): Latin Modern, sans-serif bold headings with a
%  rule under a gray chapter label, light-gray boxes, chapter frame parts
%  (Outline / Summary / Exercises / References) in small caps, exercises
%  with solutions in small type.
%
%  The first part of this file is shared verbatim with those notes, so a
%  change made there can be copied across.  The second part, below the
%  "Cast3M-specific" rule, holds what only this book needs: the gibiane
%  command boxes and the operator tables.
%======================================================================
\usepackage[margin=2.4cm,headheight=14pt]{geometry}
% Latin Modern with T1 when installed (standard TeX Live); plain Computer
% Modern otherwise, so that the file compiles on minimal installations.
\IfFileExists{lmodern.sty}{\usepackage{lmodern}\usepackage[T1]{fontenc}}{%
  \usepackage[T1]{fontenc}}
\usepackage[utf8]{inputenc}
% "english" must be named as the main language: babel otherwise takes the
% last one it processes (french) and you get French spacing and em-dash
% lists in an English book.  Use \begin{otherlanguage}{french} ... for the
% French passages quoted from the notices.
\usepackage[main=english,french]{babel}
\usepackage{amsmath,amssymb,amsthm,bm,mathtools}
\usepackage{xcolor,graphicx,enumitem,booktabs,array,multirow,tabularx}
\usepackage{verbatim,placeins}
\usepackage{tikz}
\usetikzlibrary{arrows.meta,calc,positioning,shapes.geometric,
  decorations.pathreplacing,decorations.markings,decorations.pathmorphing}
\usepackage[most]{tcolorbox}
\usepackage[font=small,labelfont={sf,bf}]{caption}
\usepackage{listings}
\usepackage{microtype}
\usepackage{etoolbox}
\usepackage{fancyhdr}
\usepackage{makeidx}
\makeindex
% The index is laid out by preamble/index.ist (sans-serif bold letter
% headings).  latexmk picks that up from .latexmkrc; by hand it is
%     makeindex -s preamble/index.ist castem
\usepackage[explicit]{titlesec}
\usepackage{hyperref}
\usepackage{url}
\usepackage{xurl}            % lets long query-string URLs break anywhere
\hypersetup{colorlinks,linkcolor=blue!50!black,citecolor=blue!50!black,
  urlcolor=blue!50!black,breaklinks=true,
  pdftitle={Introduction to Cast3M programming},
  pdfauthor={Andrei Constantinescu}}
\setlength{\emergencystretch}{3em}   % keeps long paths/URLs inside the margin
\Urlmuskip=0mu plus 1mu\relax

%---- headings: sans-serif bold, as in the RG notes ---------------------
\titleformat{\chapter}[display]
  {\sffamily\bfseries}
  {\large\color{gray!70!black}\MakeUppercase{\chaptertitlename}~\thechapter}
  {4pt}
  {\titlerule[0.6pt]\vspace{6pt}\Huge #1}
  [\vspace{2pt}]
\titleformat{name=\chapter,numberless}[display]
  {\sffamily\bfseries}{}{0pt}{\titlerule[0.6pt]\vspace{6pt}\Huge #1}
\titlespacing*{\chapter}{0pt}{-10pt}{28pt}
\titleformat{\section}{\sffamily\bfseries\Large}{\thesection}{0.8em}{#1}
\titleformat{name=\section,numberless}{\sffamily\bfseries\Large}{}{0pt}{#1}
\titleformat{\subsection}{\sffamily\bfseries\large}{\thesubsection}{0.8em}{#1}
\titleformat{name=\subsection,numberless}{\sffamily\bfseries\large}{}{0pt}{#1}
\titleformat{\subsubsection}{\sffamily\bfseries}{\thesubsubsection}{0.8em}{#1}
\titleformat{name=\subsubsection,numberless}{\sffamily\bfseries}{}{0pt}{#1}
\titleformat{\paragraph}[runin]{\sffamily\bfseries}{}{0pt}{#1}
\titlespacing*{\paragraph}{0pt}{1.6ex plus .5ex minus .2ex}{0.6em}

%---- running heads -------------------------------------------------------
\pagestyle{fancy}
\fancyhf{}
\renewcommand{\headrulewidth}{0.4pt}
\fancyhead[LE]{\sffamily\small\leftmark}
\fancyhead[RO]{\sffamily\small\rightmark}
\fancyfoot[C]{\sffamily\small\thepage}
\renewcommand{\chaptermark}[1]{\markboth{\thechapter.\ #1}{}}
\renewcommand{\sectionmark}[1]{\markright{\thesection\ \ #1}}
\fancypagestyle{plain}{\fancyhf{}\renewcommand{\headrulewidth}{0pt}%
  \fancyfoot[C]{\sffamily\small\thepage}}

%---- boxes: light gray background, bold sans-serif title -----------------
\newtcolorbox{gbox}[1]{enhanced,breakable,colback=gray!12,colframe=gray!12,
  boxrule=0pt,arc=2pt,left=8pt,right=8pt,top=6pt,bottom=6pt,
  title={#1},fonttitle=\sffamily\bfseries,coltitle=black,colbacktitle=gray!12,
  toptitle=3pt,bottomtitle=0pt,titlerule=0pt}
% untitled variant, used for take-home messages and short summaries
\newtcolorbox{gbox*}{enhanced,breakable,colback=gray!12,colframe=gray!12,
  boxrule=0pt,arc=2pt,left=8pt,right=8pt,top=6pt,bottom=6pt}

%---- theorem-like statements in the same sans-serif voice ---------------
\newtheoremstyle{sfthm}{\topsep}{\topsep}{\itshape}{}{\sffamily\bfseries}{.}{0.5em}{}
\theoremstyle{sfthm}
\newtheorem{theorem}{Theorem}[chapter]
\newtheorem{proposition}[theorem]{Proposition}
\newtheorem{lemma}[theorem]{Lemma}
\newtheorem{corollary}[theorem]{Corollary}
\newtheoremstyle{sfdef}{\topsep}{\topsep}{\normalfont}{}{\sffamily\bfseries}{.}{0.5em}{}
\theoremstyle{sfdef}
\newtheorem{definition}[theorem]{Definition}
\newtheorem{remark}[theorem]{Remark}
\renewcommand{\proofname}{\sffamily\bfseries Proof}

%---- exercises: sans-serif heading, solution in small type ---------------
\newcounter{exo}[chapter]
\renewcommand{\theexo}{\thechapter.\arabic{exo}}
\newcommand{\exercise}[1]{\par\medskip\noindent\refstepcounter{exo}%
  {\sffamily\bfseries Exercise~\theexo\quad#1.}\ }
\newcommand{\solsize}{\scriptsize}   % change to \tiny for smaller solutions
\newenvironment{solution}{\par\smallskip\begingroup\solsize\noindent
  {\sffamily\bfseries Solution.}\ }{\par\endgroup\medskip}

%---- code listings (Python), quiet gray frame in the box colour ----------
\lstdefinestyle{pystyle}{language=Python,
  basicstyle=\ttfamily\scriptsize,
  keywordstyle=\color{blue!60!black},commentstyle=\color{gray},
  stringstyle=\color{green!45!black},
  backgroundcolor=\color{gray!6},frame=none,
  breaklines=true,showstringspaces=false,tabsize=4,
  xleftmargin=0.6em,aboveskip=4pt,belowskip=4pt}
\lstset{style=pystyle}

%---- code listings (gibiane, the language of Cast3M) ---------------------
% Comments are the lines with a * in the first column (fixed-column comment
% of listings); keywords are case-insensitive, as in gibiane. Use with
% \lstinputlisting[style=gibstyle]{file.dgibi}.
\lstdefinelanguage{gibiane}{%
  sensitive=false,
  morecomment=[f][\color{gray}][0]{*},
  morestring=[b]',
  morekeywords={opti,debproc,finproc,repeter,fin,si,sinon,finsi,quitter,
    mess,table,prog,mots,lect,et,ega,dime,extr,exco,masq,vmis,sigm,bsig,
    reso,rigi,mode,mate,cara,bloq,depi,reac,pres,forc,char,evol,pasapas,
    droi,dall,cerc,coor,chan,zero,manu,dess,trac,abs,maxi,exp,log,flot,
    exis,poin,ipol,rela,elim,resu,nbno,sin,cos,inve,index,vale,chai,mot},
  keywordstyle=\color{blue!60!black}}
\lstdefinestyle{gibstyle}{language=gibiane,
  basicstyle=\ttfamily\scriptsize,
  stringstyle=\color{green!45!black},
  backgroundcolor=\color{gray!6},frame=none,
  breaklines=true,showstringspaces=false,columns=fixed,
  xleftmargin=0.6em,aboveskip=4pt,belowskip=4pt}

%---- TikZ styles shared by the diagrams ---------------------------------
\definecolor{cu}{RGB}{31,90,200}     % potential / displacement data S_u
\definecolor{cphi}{RGB}{20,150,60}   % flux / traction data S_phi, S_t
\tikzset{
  box/.style={rectangle,rounded corners=2pt,fill=gray!12,draw=none,
    minimum width=8.2cm,minimum height=1.0cm,align=center,inner sep=5pt,
    font=\small},
  smallbox/.style={rectangle,rounded corners=2pt,fill=gray!12,draw=none,
    minimum width=4.6cm,minimum height=1.25cm,align=center,inner sep=5pt,
    font=\small},
  goal/.style={fill=green!12},
  kin/.style={fill=cu!12,draw=cu!45,thin},
  stat/.style={fill=cphi!12,draw=cphi!45,thin},
  kinlab/.style={font=\small\sffamily,text=cu},
  statlab/.style={font=\small\sffamily,text=cphi},
  key/.style={fill=orange!14},
  arr/.style={-{Latex[length=2.3mm]},thick}
}

%---- framing parts of every chapter --------------------------------------
% Outline, Summary, Exercises, References: unnumbered, roman small caps,
% listed in the table of contents. \framesub is the level below (e.g. the
% parts of the exercises). A \label placed after \framepart works with
% \nameref and \pageref (not with \ref, since the part has no number).
\makeatletter
\newcommand{\framepart}[1]{%
  \par\addvspace{4.5ex plus 1ex minus .5ex}%
  \phantomsection
  \addcontentsline{toc}{section}{\protect\textsc{#1}}%
  \markright{#1}%
  \def\@currentlabelname{#1}%
  {\noindent\normalfont\Large\scshape #1\par}%
  \nobreak\vspace{0.6ex}\nobreak
  {\color{gray!55}\hrule height 0.4pt}%
  \nobreak\vspace{1.6ex plus .3ex}%
  \@afterheading}
\newcommand{\framesub}[1]{%
  \par\addvspace{3ex plus .8ex minus .3ex}%
  {\noindent\normalfont\large\scshape #1\par}%
  \nobreak\vspace{1ex plus .2ex}%
  \@afterheading}
\makeatother

%---- drafting aids (set \drafttrue/\draftfalse in castem.tex) ------------
\newif\ifdraft
\newcommand{\draftnote}[1]{\ifdraft{\sffamily\footnotesize\color{red!70!black}[#1]}\fi}

\numberwithin{equation}{chapter}
\numberwithin{figure}{chapter}
\numberwithin{table}{chapter}
\setlist{itemsep=1pt}

%======================================================================
%  Cast3M-specific: the gibiane command boxes
%
%  Every worked example in this book is a pair
%
%      \begin{castemcom}   ... explanation ...        \end{castemcom}
%      \begin{castemlines} \begin{verbatim} ... \end{verbatim} \end{castemlines}
%
%  which typesets as: explanation on the left, gibiane commands in a
%  light gray box on the right.  Keep the two environments adjacent, with
%  NO blank line between them, otherwise they end up one below the other.
%
%  The commands are captured in a save box first.  That detour is what
%  makes it legal to put a verbatim environment inside a coloured box:
%  the body of an lrbox is typeset, not read as a macro argument.
%
%  65 + 4 + 92 = 161 mm fits the 162 mm text block of the geometry above,
%  and the command box holds a 60-column gibiane line (the language itself
%  allows up to 72 characters per line).
%======================================================================

\colorlet{CodeGray}{gray!12}         % same gray as gbox, so that a command
                                     % box and a summary box match
\definecolor{Gray}{gray}{.80}        % kept for compatibility; unused

\newlength{\castemcomwidth}   \setlength{\castemcomwidth}{65.mm}
\newlength{\castemlineswidth} \setlength{\castemlineswidth}{92.mm}
\newlength{\codeboxsep}       \setlength{\codeboxsep}{5.pt}

\newsavebox{\castemsavebox}

\newenvironment{castemlines}
{\begin{lrbox}{\castemsavebox}%
 \begin{minipage}[t]{\dimexpr\castemlineswidth-2\codeboxsep\relax}%
 \begin{small}}
{\end{small}%
 \end{minipage}%
 \end{lrbox}%
 {\setlength{\fboxsep}{\codeboxsep}\colorbox{CodeGray}{\usebox{\castemsavebox}}}%
 \par\medskip\noindent}

\newenvironment{castemcom}
{\par\vspace*{3.mm}%
 \noindent
 \begin{minipage}[t]{\castemcomwidth}%
 \begin{small}}
{\end{small}%
 \end{minipage}%
 \hspace*{4.mm}}

% Full-width gray code box, for shell commands and for the Python and gmsh
% listings that have no facing comment.  Same save-box trick as castemlines.
\newsavebox{\codeboxsave}
\newenvironment{codebox}
{\par\medskip\begin{lrbox}{\codeboxsave}%
 \begin{minipage}{\dimexpr\linewidth-2\codeboxsep\relax}%
 \begin{small}}
{\end{small}%
 \end{minipage}%
 \end{lrbox}%
 \noindent{\setlength{\fboxsep}{\codeboxsep}\colorbox{CodeGray}{\usebox{\codeboxsave}}}%
 \par\medskip}

% Boxed summary table of operators:  \begin{optable}{Title} a & b \\ \end{optable}
% A gbox with the standard two columns, operator on the left in typewriter.
% The description column is 86 mm, which leaves room for the longest
% operator name in the book (31 characters) inside the box.  Note: do NOT
% use tabularx here -- inside a breakable tcolorbox it measures against
% \textwidth rather than the box content width and overflows the box.
\newenvironment{optable}[1]
{\begin{gbox}{#1}%
 \begin{tabular}{@{}>{\ttfamily}l@{\hspace{5mm}}p{86mm}@{}}}
{\end{tabular}%
 \end{gbox}}
